% from ega4-2-fr.tex
\oldpage[IV-2]{17}
Posons $B'=B\otimes_A\widehat A$~; en vertu de $1^\circ$, $B'$ est de la forme
$S^{-1}(C\otimes_A\widehat A)$, où $C$ est une $A$-algèbre de type fini et
$S$ est une partie multiplicative de $C$, donc $B'$ est un anneau noethérien.
Comme $A$ s'identifie à un sous-anneau de $\widehat A$
$(0_{\mathrm I},7.3.5)$, $A$ est intègre par $2^\circ$. L'hypothèse
$3^\circ$ entraîne alors qu'il existe un idéal premier $\mathfrak q$ de $B$
induisant l'idéal 0 de $A$ $(0_{\mathrm I},1.5.8)$, et par suite
l'homomorphisme local $A\to B/\mathfrak q$ est injectif. On peut donc se borner
à prouver la conclusion de (2.3.8) en ajoutant l'hypothèse que $B$ est un
anneau local \emph{intègre}. Appliquons (2.3.7, (ii)) à
$Y=\operatorname{Spec}(A)$, $X=\operatorname{Spec}(B)$,
$Y'=\operatorname{Spec}(\widehat A)$ et $X'=\operatorname{Spec}(B')$~; comme
le morphisme $Y'\to Y$ est plat et que $X$ (qui est intègre) domine le schéma
intègre $Y$, toute composante irréductible de $X'$ domine le schéma (intègre
par hypothèse) $Y'$. Si $y$, $x$, $y'$ sont les points fermés de $Y$, $X$ et
$Y'$ respectivement, il y a un point $x'\in X'$ (d'ailleurs unique) au-dessus
de $x$ et $y'$ (I, 3.4.9) et $\operatorname{Spec}(\mathcal O_{x'})$ domine
donc $\operatorname{Spec}(\mathcal O_{y'})$~; on a par suite un diagramme
commutatif d'homomorphismes locaux d'anneaux locaux noethériens
\[
\xymatrix{
  B=\mathcal O_x \ar[r] & \mathcal O_{x'}\\
  A=\mathcal O_y \ar[u]^{\varphi} \ar[r]_{u} &
  \mathcal O_{y'}=\widehat A \ar[u]_{v}
}
\]
tel que $u$ et $v$ soient injectifs (I, 1.2.7)~; identifiant $A$ et
$\widehat A$ à des sous-anneaux de $\mathcal O_{x'}$, et désignant par
$\mathfrak r$ l'idéal maximal de $\mathcal O_{x'}$, l'intersection des idéaux
$\mathfrak r^k\cap\widehat A$ est donc nulle $(0_{\mathrm I},7.3.5)$~;
comme $\widehat A$ est complet et que ces idéaux sont ouverts dans
$\widehat A$, cela entraîne (Bourbaki, \emph{Alg. comm.}, chap.~III, \S~2,
n\textsuperscript{o}~7, prop.~8) que la topologie de $\widehat A$ est induite
par la topologie $\mathfrak r$-préadique de $\mathcal O_{x'}$~; \emph{a
fortiori} il en est de même de la topologie de $A$
$(0_{\mathrm I},7.3.5)$. D'ailleurs on a
$\mathfrak n^k\cap A\subset\mathfrak r^k\cap A$, donc la topologie
$\mathfrak n$-préadique de $B$ induit sur $A$ une topologie plus fine que la
topologie $\mathfrak m$-préadique~; mais comme
$\mathfrak m^k\subset\mathfrak n^k\cap A$, ces deux topologies sont
identiques. C.Q.F.D.

\begin{remark}[2.3.9]
\label{IV.2.3.9-fr}
Nous verrons plus loin (7.8.3, (vii)) que pour les anneaux locaux noethériens
$A$ les plus usuels en Géométrie algébrique, l'hypothèse que $A$ est intègre
et intégralement clos implique qu'il en est de même de $\widehat A$. C'est
pourquoi, en Géométrie algébrique sur un corps de base, on énonce généralement
(2.3.8) sous l'hypothèse que $A$ est \emph{intègre et intégralement clos}.
\end{remark}

\begin{theorem}[2.3.10]
\label{IV.2.3.10-fr}
Soit $f:X\to Y$ un morphisme quasi-plat (2.3.3). Alors, pour toute partie
pro-constructible (1.9.4) $Z$ de $Y$, on a
$\overline{f^{-1}(Z)}=f^{-1}(\overline Z)$.
\end{theorem}

Comme $f$ est continue, on a
$\overline{f^{-1}(Z)}\subset f^{-1}(\overline Z)$ et tout revient à prouver
que pour tout $x\in X$ tel que $f(x)\in\overline Z$, $x$ est adhérent à
$f^{-1}(Z)$~; il est clair que la question est locale sur $Y$, donc on peut
supposer $Y$ affine. En vertu de l'hypothèse, il existe un schéma affine $Y'$
et un morphisme $g:Y'\to Y$ tels que $g(Y')=Z$ (1.9.5, (ix)). Soit $Y_1$
l'image fermée de $Y'$ par $g$ (I, 9.5.3), et soit $X_1$ le sous-préschéma
fermé $f^{-1}(Y_1)$ de $X$~; si $f_1:X_1\to Y_1$ est la restriction de $f$ à
$X_1$, on sait que $f_1$ est quasi-plat (2.3.3)~; on peut par suite remplacer
$X$, $Y$ par $X_1$, $Y_1$ respectivement, autrement dit supposer que $g$ est
\emph{dominant}. Posons alors $X'=X\times_Y Y'$, et soient $f'$ et $g'$ les
projections de $X'$ dans $Y'$ et $X$ respectivement, de sorte que l'on a le
diagramme commutatif
\[
\xymatrix{
  X \ar[d]_f & X' \ar[l]_{g'} \ar[d]^{f'} \\
  Y & Y' \ar[l]^g
}
\]
Comme $f$ est quasi-plat, $g$ quasi-compact et dominant, il résulte de (2.3.7)
(où les rôles de $f$ et $g$ sont intervertis) que $g'$ est un morphisme
dominant, ce qui démontre le théorème.

\begin{corollary}[2.3.11]
\label{IV.2.3.11-fr}
Soient $f$ un morphisme quasi-plat et quasi-compact, $F$ une partie fermée de
$X$ telle que $F=f^{-1}(f(F))$~; alors on a
$F=f^{-1}(\overline{f(F)})$.
\end{corollary}

Soient $Y'$ le sous-préschéma réduit de $X$ ayant $F$ comme espace sous-jacent
(I, 5.2.1) et soit $j:Y'\to X$ l'injection canonique~; alors $f\circ j$ est
quasi-compact (1.1.2),
